% Part: normal-modal-logic
% Chapter: filtrations
% Section: S5-decidable

\documentclass[../../../include/open-logic-section]{subfiles}

\begin{document}

\olfileid{nml}{fil}{dec}

\olsection{\Log{S5} فیصلہ کرن جوگ اے}
محدود نمونے دی خاصیت سانوں اک سوکھا طریقہ دیندی اے ایہ وکھاؤن دا کہ سانچیاں نال دِتّے امکانی منطق دے نظام \emph{فیصلہ کرن جوگ} نیں، یعنی ایہ طے کرن لئی اک قابلِ حساب طریقہ اے کہ کوئی !!a{formula} نظام وچ !!{derivable} اے یا نہیں۔

\begin{thm}
  \Log{S5} فیصلہ کرن جوگ اے۔
\end{thm}

\begin{proof}
  منّ لو $!A$ دِتّا اے تے آؤن والے !!{propositional variable}s، $!A$ وچ، $p_1$، \dots، $p_k$ وچوں نیں۔ کیوں جے ہر
  $n$ لئی $n$ جہاناں والے صرف محدود گنتی دے نمونے نیں جیہڑے
  $p_1$، \dots، $p_k$ نوں قدر دیندے نیں، اسیں \emph{نال نال} دونیں فہرستاں بنا سکدے آں۔ \Log{S5} دے سارے قضیے، کسے باقاعدہ طریقے نال ثبوت بنا کے، تے $1$، $2$، \dots جہاناں والے سارے نمونے، ایہ پرکھ کے کہ کی $!A$ ایسے کسے نمونے دے کسے جہان وچ جھوٹھا اے۔ آخر ایہناں نال نال چلدے کمّاں وچوں اک جواب دیوے گا، کیوں جے
  \olref[com][fra]{thm:generaldet} تے \olref[fmp]{cor:S5fmp} مطابق یا تاں
  $!A$، !!{derivable} اے یا کسے محدود ہمہ گیر نمونے وچ جھوٹھا اے۔
\end{proof}

اُپرلا ثبوت \Log{S5} لئی چلدا اے کیوں جے ہمہ گیر نمونیاں دیاں فلٹریشناں اپنے آپ ہمہ گیر ہوندیاں نیں۔ انعکاس تے ہر جہان توں رسائی لئی وی ایہو گل اے، پر ہور خاصیتاں لئی ہور کم چاہیدا اے۔

\end{document}
